\documentclass[a4paper,12pt]{article}

\usepackage[utf8]{inputenc}
\usepackage[top=1in,bottom=1in,left=1in,right=1in]{geometry}
\usepackage{helvet}
\usepackage[colorlinks=true,linkcolor=blue,urlcolor=blue,citecolor=blue]{hyperref}
\renewcommand{\familydefault}{\sfdefault}
\usepackage{array}
\usepackage{enumitem}
\usepackage{tocloft}
\usepackage{fancyhdr}
\usepackage{tabularx}
\usepackage[table]{xcolor}
\setlength{\parskip}{0pt}
\setlength{\parindent}{0pt}
\usepackage{microtype}
\usepackage{graphicx}
\usepackage{float}
\usepackage{longtable}
\usepackage{booktabs}

% --- Header/Footer Setup ---
\pagestyle{fancy}
\fancyhf{}
\setlength{\headheight}{15pt}
\fancyhead[L]{\small\textit{User Manual for Credit Snap}}
\fancyhead[R]{\small\textbf{Page-\thepage}}
\renewcommand{\headrulewidth}{0.6pt}

% --- Custom Colors ---
\definecolor{headergray}{RGB}{80, 80, 80}
\definecolor{linkblue}{RGB}{0, 0, 200}

% --- Helper Commands ---
\newcommand{\tocmain}[3]{%
    \noindent
    \textbf{\hyperref[#3]{\color{black} #1 \quad #2}} \dotfill \textbf{\pageref{#3}} \\[0.25cm]
}
\newcommand{\tocsub}[3]{%
    \noindent\hspace*{0.8cm}\textbf{#1\ #2} \dotfill \textbf{\pageref{#3}} \\[0.2cm]
}
\newcommand{\tocsubsub}[3]{%
    \noindent\hspace*{1.6cm}{#1\ #2} \dotfill {\pageref{#3}} \\[0.15cm]
}

% --- Gray Section Header ---
\newcommand{\sectionbar}[1]{%
    \colorbox{headergray}{%
        \makebox[\dimexpr\textwidth-2\fboxsep][c]{%
            \bfseries\textcolor{white}{\LARGE #1}%
        }%
    }%
}

\graphicspath{{./screenshots/}}

% =============================================================
\begin{document}
\pagenumbering{roman}
\thispagestyle{empty}

% ==================== COVER PAGE ====================
\begin{flushright}
    \noindent\rule{\linewidth}{4pt}\\[0.9cm]
    {\Huge\bfseries User Manual}\\[1.1cm]
    {\LARGE\bfseries for}\\[1.2cm]
    {\Huge\bfseries Credit Snap}\\[1.2cm]
    {\Large\bfseries Version 1.0}\\[1.6cm]
    {\Large\bfseries Prepared by}
\end{flushright}

\vspace{1.2cm}

\noindent
\begin{tabular*}{\textwidth}{@{} l @{\extracolsep{\fill}} r @{\hspace{1.2cm}} r @{}}
\bfseries Group 6 : &
\multicolumn{2}{r}{\bfseries\Large Group Name: \textit{Bug Dealers}} \\[1em]
\textbf{Kasavajjala Sai Shreyas}        & \textbf{240530} & \href{mailto:ksshreyas24@iitk.ac.in}{ksshreyas24@iitk.ac.in} \\
\textbf{Gade V S S R K Tejas Reddy}     & \textbf{240390} & \href{mailto:tejasreddy24@iitk.ac.in}{tejasreddy24@iitk.ac.in} \\
\textbf{Palagiri Sathish Reddy}         & \textbf{240717} & \href{mailto:psreddy24@iitk.ac.in}{psreddy24@iitk.ac.in} \\
\textbf{Kondreddy Sai Haneesh Reddy}    & \textbf{240551} & \href{mailto:haneesh24@iitk.ac.in}{haneesh24@iitk.ac.in} \\
\textbf{Yerraiahgari Ram Charan Goud}   & \textbf{241212} & \href{mailto:ramcharang24@iitk.ac.in}{ramcharang24@iitk.ac.in} \\
\textbf{Golla Lekha Harshaa}            & \textbf{240402} & \href{mailto:lekhahg24@iitk.ac.in}{lekhahg24@iitk.ac.in} \\
\textbf{Yash Raj}                       & \textbf{241202} & \href{mailto:yashr24@iitk.ac.in}{yashr24@iitk.ac.in} \\
\textbf{Korada Jayavardhan}             & \textbf{240554} & \href{mailto:jkorada24@iitk.ac.in}{jkorada24@iitk.ac.in} \\
\textbf{Boppudi Sai Chaitanya}          & \textbf{240283} & \href{mailto:boppudis24@iitk.ac.in}{boppudis24@iitk.ac.in} \\
\textbf{Panchaneni Ashwin Rao}          & \textbf{240727} & \href{mailto:ashwinr24@iitk.ac.in}{ashwinr24@iitk.ac.in} \\
\end{tabular*}

\vspace{2.5cm}

\begin{center}
\Large
\begin{tabular}{r l}
\bfseries Course:     & \bfseries CS253 \\[0.3cm]
\bfseries Mentor TA:  & \bfseries George T L \\[0.3cm]
\bfseries Date:       & \bfseries April 3, 2026
\end{tabular}
\end{center}

\vfill
\newpage
\setcounter{page}{2}

% ==================== TABLE OF CONTENTS ====================
\phantomsection
\label{toc}
\vspace*{-1cm}
\noindent
\sectionbar{CONTENTS}

\vspace{0.6cm}

\noindent\textbf{\hyperref[toc]{\color{black}CONTENTS}} \dotfill \textbf{\textcolor{black}{\pageref{toc}}} \\[0.2cm]
\noindent\textbf{\hyperref[rev]{\color{black}REVISIONS}} \dotfill \textbf{\textcolor{black}{\pageref{rev}}} \\[0.2cm]

\tocmain{1}{INTRODUCTION}{sec:introduction}
\tocsub{1.1}{PURPOSE}{sec:purpose}
\tocsub{1.2}{SOFTWARE OVERVIEW}{sec:software}
\tocsub{1.3}{PROJECT SCOPE}{sec:scope}
\tocsub{1.4}{INTENDED AUDIENCE}{sec:audience}
\tocsub{1.5}{DEFINITIONS AND ABBREVIATIONS}{sec:definitions}
\tocsub{1.6}{SYSTEM REQUIREMENTS}{sec:sysreq}

\tocmain{2}{LOGIN AND SIGNUP}{sec:login}
\tocsub{2.1}{LOGIN}{sec:login-sub}
\tocsub{2.2}{SIGN UP}{sec:signup}
\tocsub{2.3}{FORGOT PASSWORD}{sec:forgotpw}

\tocmain{3}{STUDENT MANUAL}{sec:student}
\tocsub{3.1}{HOME PAGE}{sec:stud-home}
\tocsub{3.2}{CANTEENS}{sec:stud-canteens}
\tocsub{3.3}{MENU AND ORDERING}{sec:stud-menu}
\tocsub{3.4}{VIEW DEBTS}{sec:stud-debts}
\tocsub{3.5}{HISTORY}{sec:stud-history}
\tocsub{3.6}{PROFILE}{sec:stud-profile}
\tocsub{3.7}{HELP}{sec:stud-help}

\tocmain{4}{CANTEEN OWNER MANUAL}{sec:owner}
\tocsub{4.1}{ACTIVE ORDERS DASHBOARD}{sec:owner-orders}
\tocsub{4.2}{EDIT MENU}{sec:owner-menu}
\tocsub{4.3}{ACTIVE DEBTS}{sec:owner-debts}
\tocsub{4.4}{ANALYTICS}{sec:owner-analytics}
\tocsub{4.5}{HISTORY}{sec:owner-history}
\tocsub{4.6}{PROFILE}{sec:owner-profile}

\tocmain{5}{NOTIFICATIONS}{sec:notifications}

\vspace{0.2cm}
\noindent\hyperref[app:group_log]{\color{black}\textbf{APPENDIX A --- GROUP LOG}} \dotfill \textbf{\pageref{app:group_log}} \\[0.2cm]

\newpage

% ==================== REVISIONS ====================
\phantomsection
\label{rev}
\vspace*{-1cm}
\noindent
\sectionbar{REVISIONS}

\vspace{0.6cm}

\begin{center}
\setlength{\arrayrulewidth}{1.2pt}
\renewcommand{\arraystretch}{1.5}
\begin{tabular}{|p{1.3cm}|p{5.7cm}|p{4.1cm}|p{3.2cm}|}
\hline
\rowcolor{gray!20}
\textbf{Version} & \textbf{Primary Author(s)} & \textbf{Description of Version} & \textbf{Date Completed} \\ \hline
v1.0 &
Kasavajjala Sai Shreyas \newline
Gade V S S R K Tejas Reddy \newline
Palagiri Sathish Reddy \newline
Kondreddy Sai Haneesh Reddy \newline
Yerraiahgari Ram Charan Goud \newline
Golla Lekha Harshaa \newline
Yash Raj \newline
Korada Jayavardhan \newline
Boppudi Sai Chaitanya \newline
Panchaneni Ashwin Rao
& First version of the User Manual & 03/04/26 \\ \hline
\end{tabular}
\end{center}

\newpage
\pagenumbering{arabic}

% =============================================================
%  SECTION 1 — INTRODUCTION
% =============================================================
\phantomsection\label{sec:introduction}
\sectionbar{1 INTRODUCTION}

\vspace{0.6cm}

\subsection*{1.1 PURPOSE}
\phantomsection\label{sec:purpose}

The aim of this guide is to inform the users of CreditSnap about the full functionality of the platform, detailing the operations that can be performed by each user type. By going through this guide, you will be adequately prepared to harness the entire potential of the website. The guide discusses every feature provided by the platform, describing the steps involved to perform a particular action. It also addresses common user challenges and provides clear solutions to overcome them.

\subsection*{1.2 SOFTWARE OVERVIEW}
\phantomsection\label{sec:software}

CreditSnap is a specialized MERN-stack (MongoDB, Express.js, React, Node.js) web application developed to modernize the manual operations of campus hall canteens at IIT Kanpur. Traditionally, canteens rely on paper-based registers to track student credit, which often leads to calculation errors, misplaced records, and a lack of transparency.

CreditSnap solves these issues by providing a centralized digital ecosystem with the following features:

\begin{itemize}
    \item \textbf{Dual-Interface Portal:} Custom-tailored dashboards for Students (ordering/tracking) and Canteen Owners (management/analytics).
    \item \textbf{Automated Credit Ledger:} A robust backend system that enforces credit limits and provides an immutable record of all debts and repayments.
    \item \textbf{Real-Time Communication:} Integration with Socket.io for instant order updates, ensuring students are notified the moment their meal is accepted, ready, or rejected.
    \item \textbf{Integrated Payments:} Razorpay gateway integration enables students to clear outstanding debts online (UPI, card, net banking) in real time.
    \item \textbf{Data Analytics:} An administrative suite that transforms raw transaction data into visual revenue reports and demand forecasts for canteen owners.
    \item \textbf{Institutional Security:} Secure authentication restricted to IITK credentials, ensuring the platform remains an exclusive and safe environment for campus residents.
\end{itemize}

\subsection*{1.3 PROJECT SCOPE}
\phantomsection\label{sec:scope}

The scope of CreditSnap encompasses the complete digitization of the food ordering and credit lifecycle within the IITK hall canteen system. The project aims to:

\begin{itemize}
    \item \textbf{Minimize Operational Friction:} Reduce peak-hour counter congestion by allowing students to pre-order meals digitally.
    \item \textbf{Ensure Financial Integrity:} Eliminate the ``trust gap'' between owners and students by providing real-time visibility into credit balances.
    \item \textbf{Scalability:} Support the high-traffic demands of the campus, with a system design capable of handling thousands of concurrent users and daily transactions.
    \item \textbf{Error Reduction:} Automate the calculation of daily totals and outstanding balances, which were previously prone to human error in physical ledgers.
\end{itemize}

\subsection*{1.4 INTENDED AUDIENCE}
\phantomsection\label{sec:audience}

This manual is intended for two primary user groups within the IIT Kanpur community:

\begin{enumerate}
    \item \textbf{Students:} Campus residents who wish to purchase food on credit or via pre-orders. They will use this guide to manage their profiles, track their spending, and check order statuses.
    \item \textbf{Canteen Owners:} Administrative users responsible for updating menu items, accepting or rejecting orders, managing student credit limits, and maintaining the financial health of their canteen through the digital ledger.
\end{enumerate}

No advanced technical background is required to use CreditSnap; however, a basic familiarity with web browsers and institutional email login processes is assumed.

\subsection*{1.5 DEFINITIONS AND ABBREVIATIONS}
\phantomsection\label{sec:definitions}

\vspace{0.3cm}
\renewcommand{\arraystretch}{1.4}
\setlength{\arrayrulewidth}{1pt}
\begin{tabular}{|p{3.5cm}|p{10.5cm}|}
\hline
\rowcolor{gray!20}\textbf{Term} & \textbf{Definition} \\ \hline
API & Application Programming Interface: the connection layer between the browser and the server. \\ \hline
Billing Cycle & A time period defined by a canteen owner within which outstanding student debts are expected to be cleared. \\ \hline
Debt Ceiling & The maximum credit a student may carry at a single canteen. The default value is Rs. 2500. \\ \hline
Dine-In & An order placed when the student is already present at the canteen counter. \\ \hline
HTTPS & Hypertext Transfer Protocol Secure. All data is encrypted during transmission. \\ \hline
IITK & Indian Institute of Technology Kanpur. \\ \hline
JWT & JSON Web Token: a secure short-lived digital key that confirms a user is logged in. \\ \hline
Khata & The traditional Hindi term for a credit ledger or debt book. \\ \hline
MERN & MongoDB, Express.js, React, Node.js --- the technology stack used to build CreditSnap. \\ \hline
Pre-Order & An order placed 10--40 minutes before arriving at the canteen. \\ \hline
Razorpay & The payment gateway used for online debt payments within CreditSnap. \\ \hline
SRS & Software Requirements Specification document. \\ \hline
UI & User Interface: the visual screens and controls the user interacts with. \\ \hline
\end{tabular}

\subsection*{1.6 SYSTEM REQUIREMENTS}
\phantomsection\label{sec:sysreq}

CreditSnap is a browser-based application. Users need:

\begin{itemize}
    \item A modern web browser: Google Chrome, Mozilla Firefox, Microsoft Edge, or Apple Safari.
    \item Any device: smartphone, tablet, laptop, or desktop computer.
    \item Any operating system: Windows, macOS, Linux, Android, or iOS.
    \item A stable internet connection (mobile data or Wi-Fi).
\end{itemize}

No software installation, no app download, and no disk space reservation is required.

\newpage

% =============================================================
%  SECTION 2 — LOGIN AND SIGNUP
% =============================================================
\phantomsection\label{sec:login}
\sectionbar{2 LOGIN AND SIGNUP}

\vspace{0.6cm}

When the user visits the CreditSnap website, they are presented with multiple options to navigate through the platform:\\

\textbf{Sign Up:} New users can create an account to access the platform.\\

\textbf{Login:} Existing users (students or canteen owners) can enter their registered credentials to access their accounts.\\

\textbf{Forgot Password:} If a user forgets their password, they can recover their account through the password reset process.

\subsection*{2.1 Login}
\phantomsection\label{sec:login-sub}

The Login page is designed to be intuitive, allowing users to quickly access their specific dashboards. Whether you are a student looking to order a meal or a canteen owner managing daily sales, the login process is unified.

\begin{figure}[H]
    \centering
    \includegraphics[width=1\textwidth]{Stud_Login.png}
    \label{fig:stud_login}
\end{figure}

\begin{figure}[H]
    \centering
    \includegraphics[width=1\textwidth]{Owner_Login.png}
    \label{fig:owner_login}
\end{figure}

\subsubsection*{2.1.1 Steps to Login}

\begin{enumerate}
    \item \textbf{Access the Portal:} Navigate to the CreditSnap URL. The landing page defaults to the login interface.
    \item \textbf{Select Role:} Click either the \textbf{Student} tab or the \textbf{Shopkeeper} tab depending on your account type.
    \item \textbf{Enter Credentials:} Input your registered email address (e.g., \texttt{username@iitk.ac.in} for students) and your password.
    \item \textbf{Authentication:} Click the \textbf{Login} button. The system validates your credentials against the encrypted database.
    \item \textbf{Redirection:} Upon successful validation, students are taken to the Student Home Dashboard; canteen owners are taken to the Active Orders Dashboard.
\end{enumerate}

\subsection*{2.2 Sign Up}
\phantomsection\label{sec:signup}

Registration is restricted to verified members of the IIT Kanpur community. For students, this process establishes their digital identity and initial credit profile.

\begin{figure}[H]
    \centering
    \includegraphics[width=1\textwidth]{Stud_signup.png}
    \label{fig:stud_signup}
\end{figure}

\subsubsection*{2.2.1 Creating a New Account}

To register a new student account, click the \textbf{Sign Up} link on the login page and fill in the following mandatory fields:

\begin{itemize}
    \item \textbf{Full Name:} Enter your name as it appears on your IITK Identity Card.
    \item \textbf{Institutional Email:} You must use your \texttt{@iitk.ac.in} email address. This is used for verification and official notifications.
    \item \textbf{Roll Number:} Your unique campus identifier (e.g., 240717).
    \item \textbf{Hall Number:} Select your current Hall of Residence to help canteens identify your location.
    \item \textbf{Contact Number:} A valid 10-digit mobile number for order-related updates.
    \item \textbf{Password:} Create a strong password. You will be required to enter it twice to confirm accuracy.
\end{itemize}

\subsubsection*{2.2.2 Email Verification}

After clicking \textbf{Sign Up}, the system initiates an email verification sequence:

\begin{figure}[H]
    \centering
    \includegraphics[width=1\textwidth]{Verification.png}
    \label{fig:verification}
\end{figure}

\begin{enumerate}
    \item \textbf{Verification Email:} A secure link is automatically dispatched to your IITK inbox via the Nodemailer service.
    \item \textbf{Account Activation:} You must click the link within the specified timeframe to activate your account.
    \item \textbf{Final Access:} Once verified, you will be redirected to the login page and may log in immediately.
\end{enumerate}

\subsection*{2.3 Forgot Password}
\phantomsection\label{sec:forgotpw}

If you are unable to recall your credentials, CreditSnap provides a secure password recovery workflow.

\begin{figure}[H]
    \centering
    \includegraphics[width=1\textwidth]{Forgot_Password.png}
    \label{fig:forgot_password}
\end{figure}

Follow these steps to reset your password:

\begin{enumerate}
    \item Click the \textbf{Forgot Password?} link on the login screen.
    \item Enter your registered \texttt{@iitk.ac.in} email address and click \textbf{Send Reset Link}.
    \item Check your email inbox for the password reset link sent by CreditSnap.
    \item Click the link, enter your new password, and confirm it.
    \item You will be redirected to the login page. Log in with your new credentials.
\end{enumerate}

\newpage

% =============================================================
%  SECTION 3 — STUDENT MANUAL
% =============================================================
\phantomsection\label{sec:student}
\sectionbar{3 STUDENT MANUAL}

\vspace{0.6cm}

Upon a successful login, students are directed to the primary dashboard. This page serves as the central hub for discovering campus canteens, managing personal profiles, and tracking ongoing orders. The interface is divided into a persistent navigation sidebar and a dynamic content area.

\subsection*{3.1 Home Page}
\phantomsection\label{sec:stud-home}

\begin{figure}[H]
    \centering
    \includegraphics[width=1\textwidth]{Stud_Home.png}
    \label{fig:stud_home}
\end{figure}

\subsubsection*{Sidebar Navigation}

The sidebar on the left provides quick access to the core modules of the CreditSnap platform:

\begin{itemize}
    \item \textbf{Home:} The default view showing total active debt, any important credit-limit alerts, and current active orders.
    \item \textbf{Canteens:} Displays all campus canteens with their current status (Open / Closed) and the amount you owe each one.
    \item \textbf{View Debts:} Shows a canteen-wise breakdown of your active outstanding balances, with options to pay.
    \item \textbf{History:} A complete log of all past orders and payment transactions made by the user.
    \item \textbf{Help:} Provides customer support guidelines and frequently asked questions.
\end{itemize}

\subsubsection*{Top Navigation}

On the top-right corner of the page, users can access:

\begin{itemize}
    \item \textbf{Profile:} Displays user details including Name, Roll Number, Hall, and Contact Information.
    \item \textbf{Notifications:} A bell icon that shows a badge count for unread alerts. Clicking it reveals a dropdown list of recent notifications in chronological order.
\end{itemize}

\subsubsection*{Home Dashboard Cards}

\begin{itemize}
    \item \textbf{Total Debt Card:} A prominent card at the top showing the total amount owed across all canteens, with a \textbf{Pay Now} shortcut button.
    \item \textbf{Credit Limit Alert:} A warning panel that appears when your debt at any canteen reaches 95\% or more of the allowed credit limit, naming the specific canteen.
    \item \textbf{Current Orders:} A list of all your active (non-completed) orders. Each card shows the canteen name, items ordered, quantities, total amount, payment method, order mode, and a real-time status badge (Pending / Accepted / Rejected / Ready).
\end{itemize}

\subsection*{3.2 Canteens}
\phantomsection\label{sec:stud-canteens}

Click \textbf{Canteens} in the sidebar. A list of all campus canteens is displayed, each shown as a card.

\begin{figure}[H]
    \centering
    \includegraphics[width=1\textwidth]{Stud_Canteens.png}
    \label{fig:stud_canteens}
\end{figure}

Each canteen card shows:

\begin{itemize}
    \item Canteen name and operating hours.
    \item Your current outstanding debt at that canteen.
    \item An \textbf{Open} or \textbf{Closed} status badge indicating whether orders are currently being accepted.
\end{itemize}

Click any open canteen card to proceed to its menu.

\subsection*{3.3 Menu and Ordering}
\phantomsection\label{sec:stud-menu}

After selecting a canteen, the Menu Page opens.

\begin{figure}[H]
    \centering
    \includegraphics[width=1\textwidth]{Stud_Menu.png}
    \label{fig:stud_menu}
\end{figure}

\subsubsection*{3.3.1 Browsing the Menu}

Menu items are organized into categories: \textbf{Veg}, \textbf{Non-Veg}, and \textbf{Beverages}. Each item card displays the item name, price, and category. Items marked as unavailable by the owner are grayed out and cannot be added to the cart.

\subsubsection*{3.3.2 Adding Items to Cart}

\begin{enumerate}
    \item Click \textbf{+} next to any available item to add it to the cart.
    \item Click \textbf{+} again to increase the quantity; click \textbf{---} to decrease it.
    \item The cart summary panel updates automatically with a running total.
\end{enumerate}

\subsubsection*{3.3.3 Selecting Order Mode and Payment Method}

\begin{figure}[H]
    \centering
    \includegraphics[width=1\textwidth]{Stud_Cart.png}
    \label{fig:stud_cart}
\end{figure}

Before confirming, fill in the following:

\begin{itemize}
    \item \textbf{Special Instructions:} Type any note (e.g., ``No onions'', ``Extra spicy''). This is forwarded to the canteen owner.
    \item \textbf{Order Mode:}
    \begin{itemize}
        \item \textbf{Pre-Order} --- Placing the order 10--40 minutes before arriving.
        \item \textbf{Dine-In} --- You are already at the counter or arriving immediately.
    \end{itemize}
    \item \textbf{Payment Method:}
    \begin{itemize}
        \item \textbf{Pay Now} --- Launches the Razorpay payment gateway. You pay the full amount immediately. The order is dispatched to the owner after payment is confirmed.
        \item \textbf{Add to Credit} --- The order total is added to your khata (debt) at this canteen. This option is disabled if your current debt plus the new order total would exceed your credit limit.
    \end{itemize}
\end{itemize}

Click \textbf{Confirm Order} to submit. The order appears in your Current Orders section with a \textbf{Pending} status.

\subsubsection*{3.3.4 Order Status Tracking}

The status badge on each order card updates automatically in real time:

\renewcommand{\arraystretch}{1.4}
\begin{tabular}{|p{3cm}|p{11cm}|}
\hline
\rowcolor{gray!20}\textbf{Status} & \textbf{Meaning} \\ \hline
\textbf{Pending}  & The owner has not yet acted on the order. \\ \hline
\textbf{Accepted} & The owner accepted the order. For credit orders, your debt is immediately updated. \\ \hline
\textbf{Rejected} & The owner rejected the order. No debt is recorded. You may place a new order. \\ \hline
\textbf{Ready}    & Your food is prepared and ready for pickup at the counter. \\ \hline
\end{tabular}

\vspace{0.4cm}

\subsection*{3.4 View Debts}
\phantomsection\label{sec:stud-debts}

Click \textbf{View Debts} in the sidebar. This page lists every canteen where you have an outstanding balance, with a total debt figure shown prominently at the top.

\begin{figure}[H]
    \centering
    \includegraphics[width=1\textwidth]{Stud_ViewDebts.png}
    \label{fig:stud_viewdebts}
\end{figure}

\subsubsection*{3.4.1 Paying a Debt via Razorpay}

\begin{enumerate}
    \item Find the canteen card for the debt you wish to pay.
    \item Click \textbf{Pay Now} on that card.
    \item The Razorpay payment window opens. Enter the amount (full or partial) and complete payment via UPI, card, or net banking.
    \item Once the payment is confirmed, your debt at that canteen is reduced and a record is added to your History page.
\end{enumerate}

\textit{Note for Evaluators: Razorpay is running in \textbf{Test Mode}. Do not use real banking details. Use a standard Razorpay test card to simulate a successful transaction.}

\begin{figure}[H]
    \centering
    \includegraphics[width=1\textwidth]{Razorpay.png}
    \label{fig:razorpay}
\end{figure}

\subsubsection*{3.4.2 Partial Payment}

You do not need to pay the full outstanding balance at once. Before confirming in Razorpay, modify the amount to any value less than or equal to your total debt. Your balance is reduced by however much you pay.

\subsubsection*{3.4.3 No Outstanding Debts}

If you owe nothing to any canteen, the page shows a ``No dues found'' message and the Pay Now options are not displayed.

\subsection*{3.5 History}
\phantomsection\label{sec:stud-history}

Click \textbf{History} in the sidebar. This page provides a chronological log of all your past completed orders and debt payment records.

\begin{figure}[H]
    \centering
    \includegraphics[width=1\textwidth]{Stud_History.png}
    \label{fig:stud_history}
\end{figure}

For each entry you can see:

\begin{itemize}
    \item Date and time of the transaction.
    \item Item names and quantities (for food orders).
    \item Total amount in Rs.
    \item Payment method (Online / Credit).
    \item Canteen name.
\end{itemize}

\subsection*{3.6 Profile}
\phantomsection\label{sec:stud-profile}

Click your name or avatar in the top-right header to open the Profile page.

\begin{figure}[H]
    \centering
    \includegraphics[width=1\textwidth]{Stud_Profile.png}
    \label{fig:stud_profile}
\end{figure}

The profile page displays your Full Name, Email, Hall Number, Roll Number, and Phone Number.

\subsubsection*{3.6.1 Editing Profile Information}

\begin{enumerate}
    \item Click \textbf{Edit}.
    \item Modify your \textbf{Full Name} or \textbf{Phone Number} as required.
    \item Click \textbf{Save} to apply changes.
\end{enumerate}

Note: Your email address and roll number are permanently fixed and cannot be changed. Saving will fail if the name field is left blank or the phone number is not exactly 10 digits.

\subsection*{3.7 Help}
\phantomsection\label{sec:stud-help}

Click \textbf{Help} in the sidebar. This page provides answers to frequently asked questions and guidance for common issues faced by students.

\begin{figure}[H]
    \centering
    \includegraphics[width=1\textwidth]{Stud_Help.png}
    \label{fig:stud_help}
\end{figure}

Common issues addressed include: verification email not received, Add to Credit option grayed out, order stuck on Pending, Razorpay payment errors, and profile save failures.

\newpage

% =============================================================
%  SECTION 4 — CANTEEN OWNER MANUAL
% =============================================================
\phantomsection\label{sec:owner}
\sectionbar{4 CANTEEN OWNER MANUAL}

\vspace{0.6cm}

Upon a successful login, canteen owners are directed to the Active Orders Dashboard. The interface provides a persistent sidebar and header navigation identical in layout to the student portal, but with owner-specific modules.

\subsubsection*{Sidebar Navigation}

\begin{itemize}
    \item \textbf{Home (Active Orders):} Real-time feed of all incoming student orders awaiting action.
    \item \textbf{Edit Menu:} Add, edit, or toggle the availability of menu items.
    \item \textbf{Active Debts:} View and manage outstanding balances of all students at your canteen.
    \item \textbf{Analytics:} Visual charts of revenue, order volume, and credit vs.\ cash breakdown.
    \item \textbf{History:} Log of all completed payment transactions.
    \item \textbf{Help:} Owner-specific support and FAQs.
\end{itemize}

\subsection*{4.1 Active Orders Dashboard}
\phantomsection\label{sec:owner-orders}

This is the primary working screen during operating hours. Every time a student places an order at your canteen, it appears here automatically in real time --- no page refresh required.

\begin{figure}[H]
    \centering
    \includegraphics[width=1\textwidth]{Owner_Home.png}
    \label{fig:owner_home}
\end{figure}

Each order card displays:

\begin{itemize}
    \item Student name.
    \item Items ordered and their quantities.
    \item Total cost in Rs.
    \item Payment method badge: \textbf{Added to Debt} or \textbf{Paid Online}.
    \item Order mode badge: \textbf{Pre-Order} or \textbf{Dine-In}.
    \item Two action buttons: \textbf{Accept} and \textbf{Reject}.
\end{itemize}

\subsubsection*{4.1.1 Accepting an Order}

Click \textbf{Accept} on an order card. The student receives an instant in-app notification. If the order was a credit order, the student's debt at your canteen is increased by the order total immediately. The order is removed from the Pending list.

\subsubsection*{4.1.2 Rejecting an Order}

Click \textbf{Reject} on an order card. The student receives a rejection notification. No debt is recorded regardless of the payment method. The order is removed from the active list.

\subsubsection*{4.1.3 Marking an Order as Ready}

After accepting an order and preparing the food, click \textbf{Mark as Ready}. The student receives a notification: ``Your order is ready for pickup.''

\begin{figure}[H]
    \centering
    \includegraphics[width=1\textwidth]{Owner_Order_Ready.png}
    \label{fig:owner_order_ready}
\end{figure}

\subsection*{4.2 Edit Menu}
\phantomsection\label{sec:owner-menu}

Click \textbf{Edit Menu} in the sidebar. All your current menu items are displayed in a grid of cards.

\begin{figure}[H]
    \centering
    \includegraphics[width=1\textwidth]{Owner_EditMenu.png}
    \label{fig:owner_editmenu}
\end{figure}

Each item card shows the item name, price, category, an availability toggle switch, and an \textbf{Edit} button.

\subsubsection*{4.2.1 Toggling Item Availability}

Click the toggle switch on any item:
\begin{itemize}
    \item \textbf{Toggle ON (active)} --- Item is live and visible to students.
    \item \textbf{Toggle OFF (inactive)} --- Item is hidden from the student menu immediately.
\end{itemize}

\subsubsection*{4.2.2 Editing an Existing Item}

\begin{enumerate}
    \item Click \textbf{Edit} on an item card. A popup form appears with the current details pre-filled.
    \item Modify the item name, price, or category.
    \item Click \textbf{Save}.
\end{enumerate}

\begin{figure}[H]
    \centering
    \includegraphics[width=0.75\textwidth]{Owner_EditMenu_Modal.png}
    \label{fig:owner_editmenu_modal}
\end{figure}

\subsubsection*{4.2.3 Adding a New Item}

\begin{enumerate}
    \item Click \textbf{+ Add New Item} at the top of the Edit Menu page.
    \item A form appears with blank fields.
    \item Enter the \textbf{Item Name}, \textbf{Price} (in Rs.), and select the \textbf{Category} (Veg / Non-Veg / Beverages).
    \item Set the \textbf{Availability} toggle (typically ON for a new item).
    \item Click \textbf{Save}. The item is added to your menu immediately and is visible to students.
\end{enumerate}

\begin{figure}[H]
    \centering
    \includegraphics[width=0.75\textwidth]{Owner_AddItem.png}
    \label{fig:owner_additem}
\end{figure}

\subsection*{4.3 Active Debts}
\phantomsection\label{sec:owner-debts}

Click \textbf{Active Debts} in the sidebar. This page lists every student who currently has an outstanding balance at your canteen.

\begin{figure}[H]
    \centering
    \includegraphics[width=1\textwidth]{Owner_ActiveDebts.png}
    \label{fig:owner_activedebts}
\end{figure}

Each student card shows the student's name, roll number, current outstanding debt (in Rs.), and their credit limit.

\subsubsection*{4.3.1 Searching for a Student}

Type a student's name or roll number in the search box at the top of the page. The list filters in real time. If no match is found, the message ``No matching student found'' is displayed.

\begin{figure}[H]
    \centering
    \includegraphics[width=1\textwidth]{Owner_ActiveDebts_Search.png}
    \label{fig:owner_debt_search}
\end{figure}

\subsubsection*{4.3.2 Changing a Student's Credit Limit}

\begin{enumerate}
    \item Click on a student's card to expand the detail panel.
    \item Find the \textbf{Debt Ceiling} field and enter the new limit in Rs.
    \item Click \textbf{Save}. The new limit takes effect immediately.
\end{enumerate}

Note: The new limit cannot be set below the student's current outstanding debt.

\subsubsection*{4.3.3 Setting a Default Credit Limit}

At the top of the Active Debts page there is a \textbf{Set Default Limit} field. Any new student who places their first order at your canteen automatically receives this limit.

\begin{figure}[H]
    \centering
    \includegraphics[width=1\textwidth]{Owner_DefaultLimit.png}
    \label{fig:owner_defaultlimit}
\end{figure}

\subsubsection*{4.3.4 Sending a Payment Reminder}

Click the \textbf{Notify} button on a student card. The system sends an email to that student reminding them to clear their outstanding balance.

\subsubsection*{4.3.5 Manually Recording a Cash Payment}

If a student pays you in cash, record it in the system:

\begin{enumerate}
    \item Find the student in the Active Debts list and click on their card.
    \item Enter the cash amount in the \textbf{Clear Debt / Paid Amount} field.
    \item Click \textbf{Save}. The debt is reduced and the transaction is logged in History.
\end{enumerate}

\subsection*{4.4 Analytics}
\phantomsection\label{sec:owner-analytics}

Click \textbf{Analytics} in the sidebar. This page presents visual charts summarizing your canteen's financial performance.

\begin{figure}[H]
    \centering
    \includegraphics[width=1\textwidth]{Owner_Analytics.png}
    \label{fig:owner_analytics}
\end{figure}

Charts include:

\begin{itemize}
    \item \textbf{Monthly Revenue:} Total revenue earned per month shown as a bar chart.
    \item \textbf{Orders per Month:} Number of orders processed per month.
    \item \textbf{Credit vs.\ Cash Breakdown:} Proportion of credit orders versus online payment orders.
\end{itemize}

\subsection*{4.5 History}
\phantomsection\label{sec:owner-history}

Click \textbf{History} in the sidebar. This is a chronological log of all completed transactions at your canteen --- both online payments and manually cleared cash debts.

\begin{figure}[H]
    \centering
    \includegraphics[width=1\textwidth]{Owner_History.png}
    \label{fig:owner_history}
\end{figure}

Each record displays: student name, amount paid, date and time, and payment method (Online / Cash). Records can be sorted by date, student name, or amount.

\subsection*{4.6 Profile}
\phantomsection\label{sec:owner-profile}

Click the profile icon in the header. The profile page shows your canteen name (read-only), owner name, email (read-only), and phone number.

\begin{figure}[H]
    \centering
    \includegraphics[width=1\textwidth]{Owner_Profile.png}
    \label{fig:owner_profile}
\end{figure}

Click \textbf{Edit} to modify your name or phone number, then click \textbf{Save}. Saving will fail if any required field is left blank or the phone number is not exactly 10 digits.

\newpage

% =============================================================
%  SECTION 5 — NOTIFICATIONS
% =============================================================
\phantomsection\label{sec:notifications}
\sectionbar{5 NOTIFICATIONS}

\vspace{0.6cm}

CreditSnap uses a real-time notification system to keep both students and owners informed without requiring page refreshes.

\subsection*{5.1 Student Notifications}

Students receive notifications for the following events:

\begin{itemize}
    \item Order accepted by the canteen owner.
    \item Order rejected by the canteen owner.
    \item Order is ready for pickup.
    \item Credit limit warning (when debt reaches 95\% of the allowed limit).
    \item Payment confirmation after a successful Razorpay transaction.
\end{itemize}

\subsection*{5.2 Owner Notifications}

Owners receive a notification every time a student places a new order at their canteen.

\subsection*{5.3 The Notification Bell}

The bell icon in the header shows a badge count of unread notifications. Clicking the bell opens a dropdown list of recent alerts in chronological order.

\begin{figure}[H]
    \centering
    \includegraphics[width=0.55\textwidth]{Notification_Bell.png}
    \label{fig:notification_bell}
\end{figure}

\subsection*{5.4 Email Notifications}

The following events trigger an email to the user's registered address:

\begin{itemize}
    \item \textbf{Account Verification} --- sent immediately after a new student registers.
    \item \textbf{Password Reset} --- sent when ``Forgot Password?'' is triggered.
    \item \textbf{Debt Payment Reminder} --- sent when an owner clicks the \textbf{Notify} button on a student's debt card.
\end{itemize}

\newpage

% =============================================================
%  APPENDIX A — GROUP LOG
% =============================================================
\phantomsection\label{app:group_log}
\sectionbar{APPENDIX A --- GROUP LOG}

\vspace{0.6cm}

\renewcommand{\arraystretch}{1.5}
\setlength{\arrayrulewidth}{1pt}

\begin{center}
\begin{tabular}{|p{3.5cm}|p{3.5cm}|p{7.3cm}|}
\hline
\rowcolor{gray!20}\textbf{Name} & \textbf{Roll Number} & \textbf{Contribution} \\ \hline
Kasavajjala Sai Shreyas     & 240530 & Project lead, backend, documentation coordinator \\ \hline
Gade V S S R K Tejas Reddy  & 240390 & Backend API development, Socket.io integration \\ \hline
Palagiri Sathish Reddy      & 240717 & Database design, MongoDB schema \\ \hline
Kondreddy Sai Haneesh Reddy & 240551 & Frontend --- Student portal \\ \hline
Yerraiahgari Ram Charan Goud & 241212 & Frontend --- Owner portal \\ \hline
Golla Lekha Harshaa          & 240402 & UI/UX design, CSS styling \\ \hline
Yash Raj                     & 241202 & Razorpay integration, payment flow \\ \hline
Korada Jayavardhan           & 240554 & Analytics page, charts \\ \hline
Boppudi Sai Chaitanya        & 240283 & Testing, bug fixes \\ \hline
Panchaneni Ashwin Rao        & 240727 & User Manual, SRS/SDS documentation \\ \hline
\end{tabular}
\end{center}

\end{document}
